% =====================================================================
%  Back matter -- Arabic abstract (mirrors frontmatter/abstract_en.tex)
%  Body content only: the master main.tex + preamble own all setup.
%  RTL via polyglossia \begin{Arabic} ... \end{Arabic}.
% =====================================================================
\chapter*{\textarabic{الملخص}}
\addcontentsline{toc}{chapter}{Arabic Abstract}

\begin{Arabic}
\begin{flushright}

تواجه الخدمات الإلكترونية حركة مرورٍ ترتفع وتنخفض دون سابق إنذار، بينما تأتي
الأدوات التي تحافظ على سرعتها عادةً في صورة قطعٍ منفصلة. فموازِنات الحِمل
التقليدية توجّه الطلبات وفق قواعد ثابتة لا تتكيّف مع الطلب، وأنظمة التوسّع
السحابي تضيف السعة بعد أن يكون الحِمل قد ازداد فعلاً، أمّا المراقبة ومعالجة
الأعطال فتُضاف لاحقًا. وتجميع هذه القطع معًا يمنح المشغّلين مقابض تحكّمٍ كثيرة
دون نظامٍ واحد يقرّر ما الذي ينبغي فعله. لذلك يرتفع زمن الاستجابة عند تصاعد
الطلب، وتظلّ الموارد عاطلةً عند انخفاضه، وهو ما يثقل كاهل الفِرَق الصغيرة
والمتوسطة التي لا تقوى على تشغيل منصّات تنسيق الحاويات الثقيلة.

تقدّم هذه الرسالة \textbf{SmartLoad}، وهو وسيطٌ خفيف يجمع توجيه حركة المرور،
والتنبّؤ بالطلب، وكشف الحالات الشاذّة، والتوسّع الاستباقي في حلقة قرارٍ واحدة،
مع مسارٍ احتياطيٍّ حتميٍّ عند كلّ إجراء. يراقب النظام القياسات الحيّة الواردة من
تجمّع الخوادم الخلفية، ويشغّل عدّة محرّكات قرارٍ على التوازي عبر ناقل رسائل
مشترك، ويطبّق ترتيب أولويّةٍ صارمًا: قيود السلامة أوّلًا، ثمّ التوصيات المتعلّمة
ثانيًا، ثمّ قاعدةٌ تقليديةٌ بوصفها المسار الاحتياطي المضمون. ويُسجَّل كلّ قرارٍ
آليٍّ، ويمكن للمشغّل تجاوزه يدويًّا في أيّ وقت.

بَنينا النظام في صورة حزمةٍ مُحوسَبةٍ بالحاويات، وقيّمناه على أحمال عملٍ واقعية.
وتعمل حلقة التوسّع المدفوعة بالتنبّؤ من طرفها إلى طرفها: فتحت الحِمل المتصاعد نما
تجمّع الخوادم تلقائيًّا حتى سقف سعته،
ثمّ تقلّص مجدّدًا عند انخفاض الطلب، مع جاهزية المسار الاحتياطي الحتمي في كلّ خطوة.
ويتوافق محرّك كشف الشذوذ بنسبة \textLR{\BManAgree} تقريبًا مع الأساس الحتمي المرجعي،
مع التقاطه كذلك تباينًا إضافيًّا. أمّا سياسة التوجيه المتعلّمة فتُضاهي في أدائها
أساس التناوب الدوري القياسي القويّ في وضع الظلّ، ويضمن المسار الاحتياطي المأمون
أن يظلّ التوجيه فعّالًا بالقدر نفسه الذي يحقّقه ذلك الأسلوب الراسخ.

ونخلص إلى أنّ حلقة قرارٍ متكاملةً تضع السلامة أوّلًا تمثّل أساسًا عمليًّا لإدارة
الحِمل التكيّفية: فالحلقة المغلقة تعمل من طرفها إلى طرفها، والنتائج المقيسة
تُعرَض بوضوح، ويُعدّ توسيع السلوك المتعلّم على أحمال عملٍ أشدّ تنوّعًا خطوةً
تاليةً طبيعية.

\end{flushright}
\end{Arabic}

\clearpage
